\label{sec:experiment}

\subsection{Data and training}
\label{sec:experiment-training}

\paragraph{Constructing games and histories.}
We build training data from the finite \textsc{BENAC-I} games in
Section~\ref{sec:benac-i}, varying goal dependencies, binary or linear
satisfaction, private preferences, and public priors. From these games, we
construct legal interaction histories and extract decision points together
with the information available to the acting player. We also branch on
investigation outcomes to obtain related decisions in which newly revealed
preferences can change the appropriate action. Each history is replayed
under the native game rules to verify its transitions. This construction
provides targeted examples of interpreting behavior, incorporating new
evidence, and choosing actions under different preference alignments.

\paragraph{Generating supervision.}
At each selected decision, the oracle in Section~\ref{sec:game-slices}
conditions on the player's observations and computes expected terminal
utilities for legal actions under the reference continuation. These
computations yield three training tasks: direct action selection from the
observed state and history; inference of a partner's possible and most likely
preferences; and planning given an oracle inference. Action labels retain
the acceptable actions rather than imposing a single answer when actions
tie. We exclude failed reference computations and mask planning labels that
remain ambiguous under the supplied qualitative judgment. In the paired
bank, related views and histories from the same source family share a split,
yielding 373 training and 124 validation cases. The reported checkpoints
also include an earlier inference--planning dataset and a compact joint
training set of 24 items selected using prior training performance.
Appendix~\ref{app:training} records these dataset variants and their selection
procedures.

\paragraph{Training.}
All trained variants start from Qwen3-4B-Instruct-2507~\citep{yang2025qwen3}.
The unchanged model provides the baseline. Self-play (SP) learns from each
player's terminal utility in complete games. Slice training instead uses
oracle-scored responses: \textsc{GameSlice-A} trains direct action selection,
\textsc{GameSlice-IP} trains inference and conditional planning, and
\textsc{GameSlice-Joint} combines all three tasks. These objectives update
a single model without imposing an inference-to-planning pipeline at test
time. Slice training uses grouped policy optimization with eight responses
per item and a KL penalty toward the frozen base model. We use BF16 and an
initial learning rate of $10^{-6}$. Table~\ref{tab:training-configuration}
provides the remaining settings and the realized training dose for each
verified checkpoint; the reported variants differ in both data and training
budget.

\subsection{In-domain diagnosis and adversarial games}
\label{sec:experiment-benac}

The archived post-training runs use an earlier 16-case diagnostic (v6),
distinct from the 22-case base-model diagnostic (v9) in
Section~\ref{sec:diagnose}. Table~\ref{tab:experiment-diagnosis} compares
checkpoints within that earlier protocol; its baseline scores are not
directly comparable to Table~\ref{tab:diagnose-main}.
We report exact semantic inference judgments,
planning regret when the correct judgment is supplied, native history-to-action
regret, and the paired judgment-repair effect. Regret is lower-is-better;
repair is computed only where both paired actions are valid. The three
deterministic repetitions of a structure are not independent test cases.

\begin{table}[t]
\caption{Historical post-training diagnosis under the earlier v6 protocol:
16 partner-generated histories, one per held-out structure. These are
different cases from the v9 diagnostic in Section~\ref{sec:diagnose}.
Planning regret uses the correct inference judgment.
Regret and repair are means over valid outputs; their coverage may differ.
Blank cells await a verified run under this protocol. BP-99 is a historical
inference/planning-only reference, not one of the four main arms. Valid
planning/native/repair counts are Q0 $16/15/12$, D-72 $13/16/9$,
D-148 $14/16/11$, and BP-99 $11/14/6$.}
\label{tab:experiment-diagnosis}
\centering\small
\begin{tabular}{lrrrr}
\hline
Model & Inference exact & Planning regret & Native regret & Repair $\Delta_{\mathrm{repair}}$ \\
\hline
Baseline (Q0) & 0\% & 0.844 & 0.933 & $-0.042$ \\
SP & -- & -- & -- & -- \\
O & -- & -- & -- & -- \\
D-72 & 0\% & 0.769 & 0.750 & $+0.111$ \\
D-148 & 0\% & 0.821 & 0.469 & $0.000$ \\
BP-99 (historical) & 0\% & 0.773 & 0.679 & $0.000$ \\
\hline
\end{tabular}
\end{table}

We also test full \textsc{BENAC-I} games with one focal model and two frozen
baseline partners. Initial worlds are partitioned into seen dependency
structures with new preference assignments (ID) and structurally held-out
dependencies (OOD); focal seats are rotated. We report completed episodes
and focal terminal utility conditional on completion separately, because an
invalid or truncated game has no terminal utility. The currently archived
comparison below is an earlier 48-episode, one-sample-per-scenario-seat
evaluation and does not cover the selected O and D checkpoints.

\begin{table}[t]
\caption{Historical fixed-partner adversarial games. Utility is conditional
on normal completion; completion counts are out of 24 per split. The final
common-protocol O/D comparison remains to be filled.}
\label{tab:experiment-adversarial}
\centering\small
\begin{tabular}{lrrrr}
\hline
Model & ID complete & ID utility & OOD complete & OOD utility \\
\hline
Baseline (Q0) & 22/24 & 0.288 & 22/24 & 0.068 \\
SP-69 & 20/24 & 0.525 & 18/24 & $-0.019$ \\
O & -- & -- & -- & -- \\
D & -- & -- & -- & -- \\
BP-139 (historical) & 19/24 & 0.254 & 21/24 & 0.333 \\
\hline
\end{tabular}
\end{table}

\subsection{Transfer to Complex Scenario}
\label{sec:experiment-transfer}

\paragraph{CalBench setting.}
CalBench~\citep{zou2026calbench} studies decentralized meeting scheduling:
agents see their own calendars and displacement costs, exchange messages with
other participants, and decide whether to move private errands or existing
meetings to agree on a common slot. We evaluate homogeneous teams, with the
same checkpoint controlling all four agents. Each game has eight calendar
slots and three meetings that arrive sequentially, so an early commitment can
constrain a later decision. Our frozen 24-game suite crosses four conditions
(loose availability, dense calendars, blocked errands, and a replanning
control), two errand-cost regimes (uniform or varied), and three scenario
seeds. The replanning control can require contacting a participant in an
earlier meeting who is absent from the new one. This is a small-model
sequential adaptation of CalBench, not its original aggregate benchmark score.

All models use the same cases, direct messages only, temperature zero, a
4,096-token output limit, and a 32,768-token context limit. Each meeting
allows four communication turns per participant and one decision retry; we
disable fallback and reflection. The cases have private calendar occupancy
and costs, but no future meeting or evaluator reference appears in an agent's
observation. The Qwen3-4B thinking model, MARSHAL
Generalist~\citep{yuan2026marshal}, and Social-R1-4B~\citep{wu2026socialr1}
form a separate comparison group because they use thinking-mode checkpoints;
the baseline and our trained models use Qwen3-4B-Instruct-2507.
We call the slice-trained variants \textsc{GameSlice-A} (history-to-action;
O-99), \textsc{GameSlice-IP} (inference and conditional planning; BP-99),
and \textsc{GameSlice-Joint} (all three objectives; micro24-39). The names
identify their training signals; the checkpoints also differ in exposure.

\paragraph{Evaluation.}
In our sequential suite, we count a meeting only if it remains in one common
slot for all participants in the final calendars, and count a complete game
only if all three meetings
remain valid and blocked items are preserved. Thus the first two columns of
Table~\ref{tab:experiment-calbench} have denominators 72 and 24. For a
completed game $g$, per-agent excess cost is
$(C_g-C_g^\star)/4$, where $C_g$ is realized team rescheduling cost and
$C_g^\star$ is the minimum feasible cost certified by a sequential replay;
we average this quantity over completed games only. The reference has access
to the full scenario after the fact and is not an online policy. Communication
is the native per-agent count of direct messages divided by the number of
meetings that agent helped schedule averaged
over four agents and 24 games. Burden imbalance is the native mean absolute
deviation of agents' excess burdens from their game-level mean, averaged over
the 24 games.

\begin{table}[t]
\caption{CalBench transfer on 24 frozen sequential games. Meetings count
final retained meetings; excess cost is per agent and conditional on full-game
success. Communication and imbalance average native per-agent values over all
games. Best value in each column is bold.}
\label{tab:experiment-calbench}
\centering\small
\setlength{\tabcolsep}{4pt}
\begin{tabular}{lccccc}
\hline
Model & Meetings $\uparrow$ & Full games $\uparrow$ & Cost $\downarrow$ & Comm. $\downarrow$ & Imbalance $\downarrow$ \\
\hline
Baseline & 37 & 5 & 0.100 & 5.663 & 0.443 \\
Self-play & 37 & 6 & 0.625 & 5.181 & 0.755 \\
\textsc{GameSlice-A} & 38 & 8 & 0.688 & 5.003 & 0.490 \\
\textsc{GameSlice-IP} & 44 & \textbf{10} & \textbf{0.075} & 4.760 & 0.432 \\
\textsc{GameSlice-Joint} & \textbf{47} & 9 & 0.194 & 4.714 & \textbf{0.292} \\
\hline
Qwen3-4B (thinking) & 42 & 7 & 0.571 & 4.828 & 0.557 \\
MARSHAL Generalist & 37 & 6 & 1.333 & \textbf{4.429} & 0.875 \\
Social-R1-4B & 39 & 5 & 1.100 & 4.554 & 0.901 \\
\hline
\end{tabular}
\end{table}

\textsc{GameSlice-Joint} retains the most meetings, while
\textsc{GameSlice-IP} completes the most full games. Their strengths also
differ by condition: \textsc{GameSlice-Joint} retains 17 of 18 meetings in
the loose cases and 15 of 18 in the replanning cases, compared with 14 and
10 for the unchanged baseline; \textsc{GameSlice-IP} retains 10 of 18 in
the blocked cases, compared with 6 for the baseline and 5 for
\textsc{GameSlice-Joint}. A matched initial scenario, with verbatim messages
in Appendix~\ref{app:calbench-rollout}, illustrates the
coordination behind the former result. In \texttt{replan\_varied\_s1},
\textsc{GameSlice-Joint} initially proposes slot 2, learns that it is blocked
for one partner, checks that the proposed alternative (slot 1) works for the
other, and all three participants schedule the first meeting at slot 1. When
slot 0 is later proposed for the third meeting, two participants identify its
conflict with the second meeting and the team converges on slot 2. All three
meetings remain valid. Under the same initial calendars, the baseline and
self-play teams each retain only the first meeting: in the second round a
participant reports that slot 1 is blocked, yet the resolution records
different slots.
